%! Choosing and Comparing Models in Context — Unit 2, Lesson 2.7 — Group Activity
\documentclass[10pt]{article}
\usepackage{atda-article}
\usepackage{atda-boxes}
\newcommand{\TallMath}[1]{$\displaystyle #1\rule[-1.4em]{0pt}{3.2em}$}
\setlength{\workrowsep}{6pt}

\begin{document}

\pageheader{Unit 2, Lesson 2.7}{Group Activity --- Which Family Owns This Situation?}
% NAMESTRIP: no \namedateperiod here. The student writes their name once, on the
% cover this component is stapled behind. See references/conventions.md.

\vspace{0.15in}

\begin{scenariobox}[The Scenario]{royal}
\small
\textbf{The rule for the whole activity: never name a family without naming the test.} ``Linear,
because the first differences are all $+4$'' is an answer. ``Linear, because it looks straight'' is a
guess. When the evidence is a graph instead of a table, the test is the \textbf{feature} --- a
turning point, a straight edge, a flattening end. \textbf{Everyone starts at Tier R.} When your whole
group can do Tier R without help, move on.
\end{scenariobox}

\vspace{0.10in}

\boxguard[30]
\begin{tcolorbox}[colback=white, colframe=black!40, title=\textbf{Tier R --- Remediate},
    fonttitle=\bfseries, arc=2mm, left=3mm, right=3mm, top=2mm, bottom=2mm, breakable]
\small
Two patches of duckweed are measured on the same pond every week. Both are given in square meters.

\vspace{4pt}
\renewcommand{\arraystretch}{1.3}
\begin{center}
\begin{tabular}{@{}|l|c|c|c|c|c|c|@{}}
\hline
Week $w$ & $0$ & $1$ & $2$ & $3$ & $4$ & $5$ \\ \hline
\textbf{Patch A} (m$^2$) & $6$ & $10$ & $14$ & $18$ & $22$ & $26$ \\ \hline
First differences & --- & \blank{1.0cm} & \blank{1.0cm} & \blank{1.0cm} & \blank{1.0cm} &
    \blank{1.0cm} \\ \hline
Ratios & --- & \blank{1.0cm} & \blank{1.0cm} & \blank{1.0cm} & \blank{1.0cm} & \blank{1.0cm} \\
\hline
\end{tabular}

\smallskip
\begin{tabular}{@{}|l|c|c|c|c|c|c|@{}}
\hline
Week $w$ & $0$ & $1$ & $2$ & $3$ & $4$ & $5$ \\ \hline
\textbf{Patch B} (m$^2$) & $2$ & $6$ & $18$ & $54$ & $162$ & $486$ \\ \hline
First differences & --- & \blank{1.0cm} & \blank{1.0cm} & \blank{1.0cm} & \blank{1.0cm} &
    \blank{1.0cm} \\ \hline
Ratios & --- & \blank{1.0cm} & \blank{1.0cm} & \blank{1.0cm} & \blank{1.0cm} & \blank{1.0cm} \\
\hline
\end{tabular}
\end{center}

\begin{enumerate}[leftmargin=1.6em, itemsep=4pt, topsep=4pt]
  \item Name the family for each patch, and say which row of its table decided it.

\writelines{2}

  \item Say in plain English what each patch is doing week to week. One of them \emph{adds} the same
        thing each week; the other \emph{multiplies} by the same thing.

\writelines{2}

  \item Patch A is the bigger patch in week $0$ and in week $1$. In which week does Patch B pass it?
        How much bigger is B in week $5$? Then predict both patches in week $8$.

\writelines{3}

  \item Match each patch to its graph below, and name the feature of the graph that gives it away.

\begin{center}
\begin{tabular}{@{}c@{\hspace{0.6cm}}c@{}}
\begin{tikzpicture}
\begin{axis}[
    width=5.8cm, height=3.8cm, title={\scriptsize (i)},
    xlabel={$w$ (weeks)}, ylabel={area (m$^2$)},
    xmin=0, xmax=5, ymin=0, ymax=500,
    xtick={0,1,2,3,4,5}, ytick={0,100,200,300,400,500},
    grid=both, grid style={linegray!45},
    axis lines=left, tick label style={font=\tiny}, label style={font=\tiny}]
  \addplot[royal, very thick, domain=0:5, samples=60]{2*3^x};
\end{axis}
\end{tikzpicture}
&
\begin{tikzpicture}
\begin{axis}[
    width=5.8cm, height=3.8cm, title={\scriptsize (ii)},
    xlabel={$w$ (weeks)}, ylabel={area (m$^2$)},
    xmin=0, xmax=5, ymin=0, ymax=30,
    xtick={0,1,2,3,4,5}, ytick={0,5,10,15,20,25,30},
    grid=both, grid style={linegray!45},
    axis lines=left, tick label style={font=\tiny}, label style={font=\tiny}]
  \addplot[royal, very thick] coordinates {(0,6) (5,26)};
\end{axis}
\end{tikzpicture}
\end{tabular}
\end{center}

\writelines{2}
\end{enumerate}
\end{tcolorbox}

\vspace{0.10in}

\boxguard[30]
\begin{tcolorbox}[colback=white, colframe=black!40,
    title=\textbf{Tier A --- Approaching Proficiency},
    fonttitle=\bfseries, arc=2mm, left=3mm, right=3mm, top=2mm, bottom=2mm, breakable]
\small
Three graphs, three stories. Every answer here is read off a picture.

\begin{center}
\begin{tabular}{@{}c@{\hspace{0.3cm}}c@{\hspace{0.3cm}}c@{}}
\begin{tikzpicture}
\begin{axis}[
    width=4.6cm, height=3.6cm, title={\scriptsize (a)},
    xlabel={$m$}, ylabel={$F$ (\$)},
    xmin=0, xmax=10, ymin=0, ymax=25,
    xtick={0,2,4,6,8,10}, ytick={0,5,10,15,20,25}, minor x tick num=1, minor y tick num=4,
    grid=both, grid style={linegray!45}, minor grid style={linegray!30},
    axis lines=left, tick label style={font=\tiny}, label style={font=\tiny}]
  \addplot[royal, very thick] coordinates {(0,3) (10,23)};
\end{axis}
\end{tikzpicture}
&
\begin{tikzpicture}
\begin{axis}[
    width=4.6cm, height=3.6cm, title={\scriptsize (b)},
    xlabel={$t$ (s)}, ylabel={height (ft)},
    xmin=0, xmax=4, ymin=0, ymax=70,
    xtick={0,1,2,3,4}, ytick={0,16,32,48,64}, minor x tick num=1,
    grid=both, grid style={linegray!45}, minor grid style={linegray!30},
    axis lines=left, tick label style={font=\tiny}, label style={font=\tiny}]
  \addplot[royal, very thick, domain=0:4, samples=60]{64*x-16*x^2};
\end{axis}
\end{tikzpicture}
&
\begin{tikzpicture}
\begin{axis}[
    width=4.6cm, height=3.6cm, title={\scriptsize (c)},
    xlabel={$t$ (h)}, ylabel={temp ($^\circ$C)},
    xmin=0, xmax=8, ymin=0, ymax=90,
    xtick={0,2,4,6,8}, ytick={0,20,40,60,80}, minor x tick num=1, minor y tick num=1,
    grid=both, grid style={linegray!45}, minor grid style={linegray!30},
    axis lines=left, tick label style={font=\tiny}, label style={font=\tiny}]
  \addplot[royal, very thick, domain=0:8, samples=80]{20+64*0.5^x};
\end{axis}
\end{tikzpicture}
\end{tabular}
\end{center}

\begin{enumerate}[leftmargin=1.6em, itemsep=4pt, topsep=2pt]
  \item Match each graph to its story, then name the \emph{feature of the graph} that gives the
        family away --- not just ``it looks like it.''
        \begin{enumerate}[label=(\roman*), leftmargin=1.6em, itemsep=2pt, topsep=3pt]
          \item A rideshare charges a fixed pickup fee plus a fixed amount per mile.
          \item A football is punted: it rises, turns over, and comes back down.
          \item A cup of coffee is left on a desk and cools toward room temperature.
        \end{enumerate}

\writelines{3}

  \item \textbf{Graph (a), used two ways.} A rider was charged \$17. Read off the graph how many
        miles that was. Then check it algebraically: the fare rule is $F(m)=3+2m$, so solve
        $3+2m=17$.

\begin{work}
  3 + 2m &= 17 \\
      2m &= 14 \\
       m &= 7 \text{ miles}
\end{work}

  \item \textbf{Graph (b).} Give the maximum height and the time it happens, then say when the ball
        lands. Then answer this: the model is $h(t)=64t-16t^2$, and $h(10)=-960$. Explain in one
        sentence why that is \emph{not} a mistake in the algebra.

\writelines{3}

  \item \textbf{Graph (c).} The coffee cools fast at first, then slower and slower. What temperature
        does it approach and never reach? Give the \emph{equation} of that line, and say what it is
        physically.

\writelines{2}

% Unguarded, item 5's stem sat at the foot of p2 with the figure it describes
% opening p3. \boxguard is inert inside a breakable tcolorbox, so \tcbbreak is
% the only lever. Mirrored in the other file.
\tcbbreak
  \item \textbf{The stretch --- two points are never enough.} A biologist counts $4$ cells on day
        $0$ and $8$ cells on day $1$. Both models below go exactly through those two points.

\begin{center}
\begin{tikzpicture}
\begin{axis}[
    width=8.0cm, height=3.6cm,
    xlabel={$x$ (days)}, ylabel={cells},
    xmin=0, xmax=3, ymin=0, ymax=36,
    xtick={0,1,2,3}, ytick={0,4,8,12,16,20,24,28,32,36}, minor x tick num=1,
    grid=both, grid style={linegray!45}, minor grid style={linegray!30},
    axis lines=left, tick label style={font=\tiny}, label style={font=\tiny}]
  \addplot[royal, very thick] coordinates {(0,4) (3,16)};
  \addplot[goldacc, very thick, domain=0:3, samples=60]{4*2^x};
  \addplot[royal, only marks, mark=*, mark size=2pt] coordinates {(0,4) (1,8)};
\end{axis}
\end{tikzpicture}
\end{center}

Say which curve is the linear model and which is the exponential one. Then: what \emph{one} extra
measurement would settle which model is right, and what does each model predict for it?

\writelines{3}
\end{enumerate}
\end{tcolorbox}

\vspace{0.10in}

\boxguard[30]
\begin{tcolorbox}[colback=white, colframe=black!40, title=\textbf{Tier E --- Extension},
    fonttitle=\bfseries, arc=2mm, left=3mm, right=3mm, top=2mm, bottom=2mm, breakable]
\small

\begin{enumerate}[leftmargin=1.6em, itemsep=4pt, topsep=2pt]
  \item \textbf{A claim to knock down.} A classmate writes: \emph{``The data curves upward instead of
        going straight, so the model has to be exponential.''} Name a graph from today that curves
        upward and is \textbf{not} exponential, and say what the classmate would have to check
        instead.

\writelines{2}

  \item \textbf{Two errors in one prediction.} A shop's sales, in units, for the first four months:
\smallskip

\renewcommand{\arraystretch}{1.3}
\begin{center}
\begin{tabular}{@{}|l|c|c|c|c|@{}}
\hline
Month $m$ & $0$ & $1$ & $2$ & $3$ \\ \hline
Sales & $100$ & $200$ & $400$ & $800$ \\
\hline
\end{tabular}
\end{center}

The manager looks at months $0$ and $1$, says ``sales go up by $100$ a month,'' and predicts
$100+100(12)=1300$ units in month $12$. \textbf{There are two separate errors here.} Name both, and
say what the exponential model predicts for month $12$ instead. Then say why \emph{that} number
should not be trusted either.

\writelines{4}

% Unbroken, Tier E overran p3 by a single write-line and its closing rule, so p4
% held nothing but that one line -- the worst stub in the lesson. Breaking before
% item 3 gives p4 the whole coffee item instead. Mirrored in the other file.
\tcbbreak
  \item \textbf{The coffee, done properly.} Here are the actual readings for graph (c). Room
        temperature is $20^\circ$C.
\smallskip

\renewcommand{\arraystretch}{1.3}
\begin{center}
\begin{tabular}{@{}|l|c|c|c|c|c|c|c|@{}}
\hline
$t$ (hours) & $0$ & $1$ & $2$ & $3$ & $4$ & $5$ & $6$ \\ \hline
$C$ (temperature, $^\circ$C) & $84$ & $52$ & $36$ & $28$ & $24$ & $22$ & $21$ \\ \hline
Ratios of $C$ & --- & \blank{1.0cm} & \blank{1.0cm} & \blank{1.0cm} & \blank{1.0cm} &
    \blank{1.0cm} & \blank{1.0cm} \\ \hline
Gap above $20^\circ$ & \blank{1.0cm} & \blank{1.0cm} & \blank{1.0cm} & \blank{1.0cm} &
    \blank{1.0cm} & \blank{1.0cm} & \blank{1.0cm} \\ \hline
Ratios of the gap & --- & \blank{1.0cm} & \blank{1.0cm} & \blank{1.0cm} & \blank{1.0cm} &
    \blank{1.0cm} & \blank{1.0cm} \\
\hline
\end{tabular}
\end{center}

\begin{enumerate}[label=(\alph*), leftmargin=1.6em, itemsep=3pt, topsep=4pt]
  \item The ratios of $C$ are not constant --- so is the ratio test telling you the coffee is
        \emph{not} exponential? Look at the last row before you answer.

\writelines{2}

  \item The gap between the coffee and the room is what halves. Write down what transformation
        (Lesson 2.2) turns the parent exponential into this graph, and give the equation of the
        horizontal asymptote.

\writelines{2}
\end{enumerate}
\end{enumerate}
\end{tcolorbox}

\end{document}
